
% CREACIÓN DEL DOCUMENTO
\documentclass[
	spanish, % Idioma: spanish, english, etc.
	letterpaper, oneside
]{article}

% INFORMACIÓN DEL DOCUMENTO
\def\documenttitle {Cuadro comparativo de criterios de la redacción académica}
\def\documentsubtitle {Actividad 2 - Redacción en contextos virtuales}
\def\documentsubject { Actividad Semana 3}

\def\documentauthor {Martín Jonathan de la Cruz}
\def\coursename {Redacción en contextos virtuales}
\def\coursecode {017}

\def\universityname {Universidad Abierta y a Distancia de México}
\def\universityfaculty {Licenciatura en Derecho}
\def\universitydepartment {Redacción en contextos virtuales}
\def\universitydepartmentimage {departamentos/UnADM}
\def\universitydepartmentimagecfg {height=1.57cm}
\def\universitylocation {Roma Norte, Ciudad de México}

% INTEGRANTES, PROFESORES Y FECHAS
\def\authortable {
	\begin{tabular}{ll}
		\textbf{Alumno:}
		& \begin{tabular}[t]{l}
			 Martín Jonathan de la Cruz \\	
		\end{tabular} \\
		Matrícula: &   ES2611202040 \\
		
		\textit{Profesor:}
		& \begin{tabular}[t]{l}
			EAna \\ Arceo Luna
		\end{tabular} \\
		& \\
		\multicolumn{2}{l}{Fecha de realización: \today} \\
		% \multicolumn{2}{l}{Fecha de entrega: \today} \\
		\multicolumn{2}{l}{\universitylocation}
	\end{tabular}
}

% IMPORTACIÓN DEL TEMPLATE
\input{template}

% CITAS EN FORMATO AUTOR-AÑO (APA)
\setcitestyle{authoryear,open={(},close={)}}

% INICIO DE PÁGINAS
\begin{document}

% PORTADA
\templatePortrait

% CONFIGURACIÓN DE PÁGINA Y ENCABEZADOS
\templatePagecfg

% RESUMEN O ABSTRACT
\begin{abstractd}
	\textbf{Objetivo:} Identificar y aplicar los criterios básicos de la redacción académica --claridad, coherencia, cohesión y formalidad-- para valorar su función en la producción de textos dentro de entornos educativos virtuales.

	\textbf{Actividad realizada:} Se revisó la bibliografía indicada y se analizaron dos textos sobre educación en línea. Con base en la técnica de cuadro comparativo, se contrastaron ambos textos según los criterios de claridad, coherencia, cohesión, formalidad y objetividad. Finalmente, se integró una lectura crítica sobre la importancia de la redacción académica en la educación a distancia.
\end{abstractd}

% TABLA DE CONTENIDOS - ÍNDICE
\templateIndex

% CONFIGURACIONES FINALES
\templateFinalcfg

% ======================= INICIO DEL DOCUMENTO =======================

\section{Introducción}

La redacción académica no es una exigencia decorativa: es un sistema de control del sentido. En la educación a distancia, donde la mayor parte de la interacción ocurre por escrito, una idea mal organizada puede convertirse en confusión, retraso, baja participación o interpretación equivocada de las instrucciones. Por eso, escribir con claridad, coherencia, cohesión y formalidad no solo mejora la presentación de un texto; también ordena el aprendizaje y reduce el costo comunicativo de estudiar en plataformas virtuales.

El análisis se realiza mediante un cuadro comparativo, técnica didáctica reconocida por la Universidad Abierta y a Distancia de México dentro de sus recursos de enseñanza y aprendizaje \citep{unadm100TecnicasDidacticas2023}. Esta técnica permite observar diferencias de estructura, tono y función entre dos textos. La bibliografía revisada sostiene que la escritura académica exige planificación, redacción y revisión \citep{nunezcortesEscrituraAcademicaTeoria2015}; además, requiere claridad para evitar ambigüedades \citep{ayalaTallerLecturaRedaccion2010} y debe entenderse como una forma de comunicación que articula vocabulario, cohesión y coherencia \citep{vitoriaEscrituraAcademicaFormacion2019}. El problema operativo, entonces, no es decidir cuál texto ``suena mejor'', sino identificar cuál funciona mejor para producir comprensión académica.

\section{Desarrollo}

Se analizarán dos textos que abordan la experiencia de estudiar en línea.

\textbf{Texto 1:} ``La neta siento que estudiar en línea está medio complicado porque luego los profes dejan muchas tareas y ni explican bien, o sea sí entiendo que es responsabilidad de uno pero pues también deberían apoyar más. A veces entro a la plataforma y no entiendo ni qué hacer primero, todo está revuelto y pues te estresas. Igual creo que sí se aprende pero depende mucho de cada quien y de cómo le eche ganas.''

\textbf{Texto 2:} ``La educación en línea representa una modalidad que exige altos niveles de autonomía por parte del estudiante. No obstante, uno de los principales desafíos que enfrentan los alumnos es la organización de las actividades y la claridad en las instrucciones proporcionadas por el docente. Cuando estas no son suficientemente precisas, pueden generar confusión y afectar el proceso de aprendizaje. 

En este sentido, es fundamental que exista una comunicación clara y estructurada en las plataformas educativas, así como un acompañamiento docente que favorezca la comprensión de las tareas. De esta manera, se promueve un aprendizaje más efectivo y significativo en los entornos virtuales.'' 

\subsection{Cuadro comparativo}

El Texto 1 expresa una experiencia personal sobre estudiar en línea, pero lo hace desde un registro coloquial y con baja organización argumentativa. El Texto 2 aborda el mismo problema, aunque lo transforma en una explicación académica: identifica una causa, describe una consecuencia y propone una condición de mejora.

\clearpage
\pagestyle{empty}
\begin{landscape}
\begingroup
\scriptsize
\setlength{\tabcolsep}{4pt}
\renewcommand{\arraystretch}{1.35}
\begin{longtable}{@{}p{0.14\linewidth}p{0.27\linewidth}p{0.27\linewidth}p{0.28\linewidth}@{}}
\caption{Comparación de criterios de redacción académica en los textos analizados}
\label{tab:comparacion-redaccion-academica}\\
\toprule
\textbf{Criterio} & \textbf{Texto 1} & \textbf{Texto 2} & \textbf{Diferencia principal} \\
\midrule
\endfirsthead
\toprule
\textbf{Criterio} & \textbf{Texto 1} & \textbf{Texto 2} & \textbf{Diferencia principal} \\
\midrule
\endhead
\midrule
\multicolumn{4}{r}{Continúa en la siguiente página} \\
\endfoot
\bottomrule
\endlastfoot

\textbf{Claridad} &
Comunica inconformidad, pero no delimita con precisión el problema. Mezcla exceso de tareas, falta de explicación docente, desorden de plataforma y esfuerzo individual sin jerarquizar qué aspecto causa la dificultad. El lector entiende el malestar general, pero debe inferir el punto central. &
Define con mayor precisión el tema: la educación en línea exige autonomía y puede verse afectada por instrucciones poco claras. Presenta una relación causa-efecto visible: si las indicaciones no son precisas, se produce confusión y se afecta el aprendizaje. &
El Texto 1 depende de la experiencia emocional; el Texto 2 controla el sentido mediante conceptos y relaciones explícitas. La claridad no elimina la postura, pero obliga a convertirla en información comprensible. \\
\midrule

\textbf{Coherencia} &
La progresión de ideas es débil. El texto pasa de quejas sobre docentes a problemas de plataforma y luego a una afirmación general sobre ``echarle ganas''. No existe una tesis estable ni un cierre que ordene las ideas anteriores. &
Mantiene una línea argumentativa: primero plantea la exigencia de autonomía, después identifica un desafío y finalmente propone comunicación estructurada y acompañamiento docente. Cada parte cumple una función dentro del argumento. &
El Texto 1 funciona por asociación de ideas; el Texto 2 funciona por estructura. La coherencia aparece cuando cada oración responde a una dirección común y no solo a la acumulación de impresiones. \\
\midrule

\textbf{Cohesión} &
Usa conectores propios del habla cotidiana, como ``porque'', ``pues'', ``o sea'' e ``igual'', que enlazan frases de manera informal, pero no siempre establecen relaciones lógicas precisas. La repetición de expresiones debilita la continuidad textual. &
Emplea conectores y referencias internas con mayor control: ``no obstante'', ``cuando'', ``en este sentido'' y ``de esta manera''. Estos recursos vinculan problema, consecuencia y propuesta, lo que permite seguir el desarrollo sin rupturas. &
El Texto 1 une frases; el Texto 2 articula ideas. La cohesión no consiste en agregar conectores, sino en usarlos para marcar causa, contraste, continuidad y consecuencia. \\
\midrule

\textbf{Formalidad} &
Predomina un registro coloquial: ``la neta'', ``medio complicado'', ``profes'', ``pues'' y ``echarle ganas''. Ese tono puede funcionar en una conversación informal, pero pierde fuerza en un contexto académico porque reduce precisión y autoridad. &
Utiliza un registro académico e impersonal: ``modalidad'', ``autonomía'', ``desafíos'', ``organización de las actividades'', ``proceso de aprendizaje'' y ``entornos virtuales''. El lenguaje se ajusta mejor a una actividad universitaria. &
El Texto 1 traslada códigos de conversación cotidiana al aula virtual; el Texto 2 adapta el discurso al contexto universitario. La formalidad opera como criterio de adecuación, no como simple adorno. \\
\midrule

\textbf{Objetividad} &
Se apoya en una percepción subjetiva y generaliza sin evidencia. Expresiones como ``ni explican bien'' o ``todo está revuelto'' señalan un problema posible, pero no ofrecen criterios verificables para analizarlo. &
Formula el problema de manera más objetiva. No acusa directamente a personas, sino que describe condiciones del proceso educativo: autonomía, organización, claridad de instrucciones y acompañamiento docente. &
El Texto 1 denuncia una molestia; el Texto 2 transforma la molestia en objeto de análisis. Esa conversión es clave en la redacción académica, porque permite discutir el problema sin reducirlo a reacción personal. \\
\end{longtable}
\endgroup
\end{landscape}

\pagestyle{fancy}

\subsection{Análisis del contraste}

El contraste muestra que la redacción académica cumple una función estratégica: convierte experiencias dispersas en problemas analizables. El Texto 1 no es inútil, porque revela una tensión real de la educación en línea: la carga de actividades, la necesidad de acompañamiento y la posible confusión dentro de la plataforma. Sin embargo, su forma le resta fuerza, porque el registro coloquial y la falta de estructura hacen que el lector se concentre más en el tono que en el problema. El Texto 2 gana eficacia porque organiza la información, evita ataques personales y plantea una salida institucional: comunicación clara, instrucciones estructuradas y acompañamiento docente. En términos académicos, gana quien logra que su escritura sea verificable, ordenada y útil para tomar decisiones.

\section{Conclusión}

La redacción académica es decisiva en la educación a distancia porque sustituye buena parte de la explicación presencial. Cuando el texto es claro, el estudiante sabe qué debe hacer; cuando es coherente, entiende por qué debe hacerlo; cuando es cohesivo, sigue el razonamiento sin perderse; y cuando es formal, reconoce el contexto universitario. Mi postura es que escribir bien no es solo una habilidad escolar, sino una herramienta de autonomía. En entornos virtuales, quien domina la escritura controla mejor su aprendizaje, participa con mayor autoridad y reduce errores de interpretación que afectan el desempeño académico.

\bibliography{RedaccionContextosVirtuales}

% FIN DEL DOCUMENTO
\end{document}
